% !PRINTED 176
\weberpage{176}
%<<CONT
are. Consequently (having regard to \S\,45) one of the following
relations holds:
\[
\lambda^2 = \varkappa^2(\omega), \quad \varkappa^2\!\left(-\frac{1}{\omega}\right), \quad \varkappa^2(\omega+1), \quad \varkappa^2\!\left(\frac{\omega-1}{\omega}\right),
\]
\[
\varkappa^2\!\left(-\frac{1}{\omega+1}\right), \quad \varkappa^2\!\left(\frac{\omega}{1-\omega}\right),
\]
whereby the correctness of the second theorem follows from the first.

The variable $\omega$ is here always restricted to a positive
imaginary part, and if we call two numbers $\omega$, $\omega'$ that are
connected by an equation (11) \gesperrt{equivalent}, then all numbers
equivalent to $\omega$ have a positive imaginary part.
We can also state our theorem 2. in this way:

3. \gesperrt{The function $j(\omega)$ has one and the same value for all equivalent values $\omega$
and only for these}$^{1)}$.

To every value of $\omega$ there belongs a definite value of $j(\omega)$,
and to every value of $j(\omega)$ a value of $\omega$ and the totality
 of the equivalent values $\omega'$. If therefore a one-valued
function $\Phi(\omega)$ of $\omega$ satisfies the condition:
\[
\Phi\!\left(\frac{\gamma+\delta\omega}{\alpha+\beta\omega}\right) = \Phi(\omega),
\]
then $\Phi(\omega)$ is a one-valued function of $j(\omega)$, and if we
therefore assume that $\Phi(\omega)$ becomes infinite only for a finite
number of values $j(\omega)$ and only in finite order, it follows
that $\Phi(\omega)$ is a \gesperrt{rational function of $j(\omega)$}.

The assumption always holds, for example, when $\Phi(\omega)$
stands in an \gesperrt{algebraic connection} with $j(\omega)$.

\begin{center}
\textbf{\S\ 54. \gesperrt{The Modular Functions.}}
\end{center}

Let $\psi(\omega)$ be a single-valued function of $\omega$. If one applies
to $\omega$ a linear transformation $S = \begin{pmatrix} \alpha, & \beta \\ \gamma, & \delta \end{pmatrix}$, then $\psi(\omega)$
passes into another function
\weq{1}{\psi\!\left(\frac{\gamma+\delta\omega}{\alpha+\beta\omega}\right)}
which we shall denote by $\psi\,|\,S$.

\begingroup
\renewcommand{\thefootnote}{1)}
\footnotetext[1]{This theorem was first proved by \gesperrt{Dedekind} \gesperrt{(Crelle, Bd.\,83)}.
\gesperrt{Dedekind} accordingly calls the function $\mathrm{val}(\omega) = 3^{-3}2^{-6}j(\omega)$ the \gesperrt{valence}
of $\omega$.}
\endgroup
